% from main.tex
%%%%%%%%%%  scan idx 529  |  folio 516  %%%%%%%%%%
Le
groupe $G$ opère sur $E_0$, sur $X_0$ muni des faisceaux $L_0^{(i)}$ et sur $Y_0$. Si on fait le produit de ces objets avec $P$, sous $G$, on obtient le faisceau $E$ associé à $P$, le schéma $X=D(E)=P/B$ des drapeaux de $E$, et le schéma $Y=P/U$. On notera les faisceaux canoniques sur $X$ par les mêmes lettres que leurs homologues sur $X_0$. On a donc le diagramme :
\[
\begin{array}{ccccc}
&& Y=P/U &&\\[-0.2em]
& \nearrow && \searrow &\\[-0.2em]
P &&&& X=P/B=D(E)\\[-0.2em]
& \searrow && \swarrow &\\[-0.2em]
&& S &&
\end{array}
\]
On remarque que $P$ est un torseur de base $Y$, de groupe structural $U_Y$, et que $Y$ est un torseur de base $X$, de groupe structural $T_X$, défini par la $\mathcal O_X$-Algèbre :
\[
A=\bigoplus_{k_i\in\mathbb Z}(F^{(1)})^{\otimes k_1}\otimes\cdots\otimes(F^{(n)})^{\otimes k_n}.
\]

L'homomorphisme $K_\bullet(S)\to K_\bullet(P)$ se factorise donc en :
\[
K_\bullet(S)\longrightarrow K_\bullet(X)\longrightarrow K_\bullet(Y)\longrightarrow K_\bullet(P).
\]
D'après 3.1, on a $K_\bullet(X)=K_\bullet(S)\otimes_{K^\circ(S)}K^\circ(X)$. Soit $\xi_i$ la classe de $F^{(i)}$ dans $K^\circ(X)$. Alors d'après 2.2, $K_\bullet(Y)$ est le quotient de $K_\bullet(X)$ par le sous-module engendré par les $(\xi_i-1)K_\bullet(X)$. Si $I$ est l'idéal de $K^\circ(X)$ engendré par les $\xi_i-1$, on a donc :
\[
K_\bullet(Y)=K_\bullet(S)\otimes_{K^\circ(S)}(K^\circ(X)/I).
\]
